\section{Описание алгоритма}

Описанный далее алгоритм основан на алгоритме <<UNIFY-1>>, представленном в работе Мартелли и Росси~\cite{martelli1984efficient}. По сравнению с ним, наш алгоритм специализирован для обработки одного уравнения за раз, что зачастую и требуется на практике, вместо обработки заранее заданного множества уравнений. Это позволило упростить описание структур данных и самого алгоритма, приблизив их к практической реализации.

\subsection{Система уравнений}

Несмотря на то, что понятие системы уравнений обладает общепринятым определением, в данной работе мы так называем объект, получаемый в результате работы нашего алгоритма.

Для некоторого множества $A$ введём обозначение $A^? = \{ \bot \} \cup A$, которое мы называем <<опциональный $A$>>.

\begin{definition}[Уравнение]
    Назовём \textit{уравнением} $e \in \E$ упорядоченную пару из:
    \begin{enumerate}
        \item множества переменных <<левая сторона>> с отмеченной точкой <<представляющая переменная>>;
        \item опционального (конечного) терма <<правая сторона>>, не являющегося переменной.
    \end{enumerate}

    Для удобства мы вводим проекции $\lhs : \E \rightarrow 2^\V$, $\repr : \E \rightarrow \V$ и $\rhs : \E \rightarrow T^?$ для левой стороны, представляющей переменной и правой стороны уравнения соответственно. Также, мы используем нотацию $\underline{x}, y, z \equiv t$ для обозначения уравнения, в котором $\lhs(\underline{x}, y, z \equiv t) = \{ x, y, z \}$, $\repr(\underline{x}, y, z \equiv t) = x$ и $\rhs(\underline{x}, y, z \equiv t) = t$.
\end{definition}

\begin{definition}[Система уравнений]
    Назовём \textit{системой уравнений} $\xi \in \Xi$ конечное множество уравнений с попарно непересекающимися левыми сторонами. Для удобства мы будем обозначать пополнение системы уравнений новым уравнением $e$ как $\xi, e$:
    \[\{ e \} \cup \{ e' \in \xi \mid \lhs(e) \cap \lhs(e') = \emptyset \}.\]
\end{definition}

Данное определение уравнения стоит рассматривать как частный случай системы уравнений в общепринятом смысле. Например:
\begin{equation*}
    \begin{array}{c|c}
        \text{Общепринятая система уравнений} & \text{Наше уравнение} \\[2pt]
        \left\{\begin{array}{l@{{}={}}l}
            x & f\,(z) \\
            y & x \\
            z & x \\
        \end{array}\right. & \underline{x}, y, z \equiv f\,(z) \\
    \end{array}
\end{equation*}

В этом смысле тривиальные уравнения вида $\underline{x} \equiv \bot$ <<ничего не уравнивают>>, поэтому мы дополнительно считаем, что никакая система уравнений $\xi \in \Xi$ не может содержать таких уравнений:
\[\forall x \in \V, \xi \in \Xi. ~ \underline{x} \equiv \bot \notin \xi.\]

Мы определяем две операции на системах уравнений. Основная операция $\walk : \Xi \times \V \rightarrow \E \times T$ по данной системе уравнений и переменной находит соответствующее уравнение:
\begin{equation*}
    \walk(\xi, x) = \begin{cases}
        (e, \rhs(e)) & \text{если $e \in \xi, x \in \lhs(e)$ и $\rhs(e) \neq \bot$}, \\
        (e, \repr(e)) & \text{если $e \in \xi, x \in \lhs(e)$ и $\rhs(e) = \bot$}, \\
        (\underline{x} \equiv \bot, x) & \text{иначе}. \\
    \end{cases}
\end{equation*}

С помощью этой операции можно определить отображение $\llbracket \cdot \rrbracket : \Xi \rightarrow \Sigma$, задающее интерпретацию системы уравнений как подстановки. Результат применения $t \llbracket \xi \rrbracket$ задаётся следующим псевдокодом:
\begin{algorithmic}
    \Function{ApplyEqSys}{$\xi \in \Xi$, $t \in T^\infty$}
        \If{$x \gets t \in \V$}
            \State $(\_, t) \gets \walk(\xi, x)$
        \EndIf
        \If{$x' \gets t \in \V$}
            \State \Return $x'$
            \Comment{неподвижная точка}
        \EndIf
        \If{$f^n\,(t_1, \dots, t_n) \gets t$}
            \State \Return $f^n\,(\Call{ApplyEqSys}{\xi, t_1}, \dots, \Call{ApplyEqSys}{\xi, t_n})$
        \EndIf
    \EndFunction
\end{algorithmic}

\begin{theorem}
    Для любых $\xi \in \Xi$ и $t \in T^R$ верно, что $t \llbracket \xi \rrbracket \in T^R$.
\end{theorem}
\begin{proof}
    Достаточно заметить, что уравнений в системе конечное число и их правые стороны также конечны. Более подробно см. \texttt{wf\_eqsys\_to\_subst\_rational} и \texttt{inf\_subst\_apply\_rational} в реализации на \Rocq.
\end{proof}

Вторая вводимая операция $\union : \Xi \times \V \times \V \rightarrow \Xi \times (T \times T)^?$ преобразует данную систему уравнений так, чтобы данные переменные лежали в одном уравнении:
\begin{algorithmic}[1]
    \Function{UnionEqSys}{$\xi \in \Xi$, $x, y \in \V$}
        \State $(e_x, \_) \gets \walk(\xi, x)$
        \State $(e_y, \_) \gets \walk(\xi, y)$
        \If{$\repr(e_x) = \repr(e_y)$}
            \State \Return $(\xi, \bot)$
            \Comment{уже лежат в одном уравнении или равны}
        \EndIf
        \If{$\rhs(e_x) = \bot$ \textbf{and} $\rhs(e_y) = \bot$}
            \State \Return $([\xi, \lhs(e_x) \cup \lhs(e_y) \equiv \bot], \bot)$
        \EndIf
        \If{$\rhs(e_x) = \bot$}
            \State \Return $([\xi, \lhs(e_x) \cup \lhs(e_y) \equiv \rhs(e_y)], \bot)$
        \EndIf
        \If{$\rhs(e_y) = \bot$}
            \State \Return $([\xi, \lhs(e_x) \cup \lhs(e_y) \equiv \rhs(e_x)], \bot)$
        \EndIf
        \State \Return $([\xi, \lhs(e_x) \cup \lhs(e_y) \equiv \rhs(e_x)],(\rhs(e_x), \rhs(e_y)))$ \\
        \Comment{конфликт правых сторон уравнений}
    \EndFunction
\end{algorithmic}

В приведённом псевдокоде, когда добавляется новое уравнение на строках 6-12, мы можем выбрать в качестве представляющей переменной как $\repr(e_x)$, так и $\repr(e_y)$ -- это не влияет на теоретические свойства алгоритма. То же верно и про выбор правой стороны добавляемого уравнения на строке 12: это может быть как $\rhs(e_x)$, так и $\rhs(e_y)$. Как показывает псевдокод, второе возвращаемое значение сигнализирует о конфликте правых сторон. Квадратные скобки используются для группировки и не несут специального смысла.


\subsection{Общая часть термов}

Одной из важных частей алгоритма является операция выделения <<общей части>>~\cite{martelli1984efficient} двух конечных термов, $\commonpart : T \times T \rightarrow T^?$:
\begin{equation*}
    \begin{split}
        \commonpart(x, \_) &= x \\
        \commonpart(\_, y) &= y \\
        \commonpart(f^n\,(..., l_i, ...), f^n\,(..., r_i, ...)) &= f^n\,(..., \commonpart(l_i, r_i), ...) \\
    \end{split}
\end{equation*}

Мы задаём поведение этой функции в стиле индуктивного предиката. Такое определение удобно, потому что оставляет некоторую свободу в выборе результата, когда функция применяется к двум переменным -- можно выбрать любую из них. В случаях, когда никакое из правил не подходит, полагается значение $\bot$.


\subsection{Алгоритм}

Алгоритм задаётся рекурсивной функцией $\unify : \Xi \times T \times T \rightarrow \Xi^?$:
\begin{algorithmic}[1]
    \Function{RationalUnify}{$\xi \in \Xi$, $t_1, t_2 \in T$}
        \If{$x \gets t_1 \in \V$ \textbf{and} $y \gets t_2 \in \V$}
            \State $(\xi, ts) \gets \union(\xi, x, y)$
            \If{$(t_1, t_2) \gets ts$}
                \State \Return $\Call{RationalUnify}{\xi, t_1, t_2}$
                \Comment{разрешение конфликта}
            \EndIf
            \State \Return $\xi$
        \EndIf
        \If{$x \gets t_1 \in \V$}
            \State $(e, \_) \gets \walk(\xi, x)$
            \If{$\rhs(e) = \bot$}
                \State \Return $\xi, \lhs(e) \equiv t_2$
            \EndIf
            \State $t \gets \commonpart(\rhs(e), t_2)$
            \If{$t = \bot$}
                \State {\bf failure}
            \EndIf
            \State $\xi \gets \xi, \lhs(e) \equiv t$
            \Comment{упрощение правой стороны уравнения}
            \State \Return \Call{RationalUnify}{$\xi, \rhs(e), t_2$}
        \EndIf
        \If{$y \gets t_2 \in \V$}
            \State \Return \Call{RationalUnify}{$\xi, y, t_1$}
        \EndIf
        \If{$f^n\,(l_1, ..., l_n) \gets t_1$ \textbf{and} $g^m\,(r_1, ..., r_m) \gets t_2$}
            \If{$f^n \neq g^m$}
                \State \bf failure
            \EndIf
            \For{$i = 1, ..., n$}
                \State $\xi \gets \Call{RationalUnify}{\xi, l_i, r_i}$
            \EndFor
            \State \Return $\xi$
        \EndIf
    \EndFunction
\end{algorithmic}

\begin{theorem}[Корректность]
    Алгоритм \textsc{RationalUnify} корректен.
\end{theorem}
\begin{proof}
    Точная формулировка корректности и доказательство (индукция по дереву вывода) выходят за рамки данной статьи, но приведены в реализации на \Rocq под названиями \texttt{rational\_unification\_correct} и \texttt{rational\_unification\_complete}.
\end{proof}

\begin{theorem}[Завершаемость]
    Для любых $\xi \in \Xi$ и $t_1, t_2 \in T$ алгоритм \textsc{RationalUnify} завершается с некоторым результатом.
\end{theorem}
\begin{proof}
    Фундированная (well-founded) индукция по <<размеру>> входных параметров алгоритма. См. \texttt{rational\_unification\_exists} в реализации на \Rocq.
\end{proof}
